\scanpage{152}
\mspage{26}
% archive sheet number: 151

de la forme
\[
s'=\deleted{\unsure{s^{\pi}}}s^t,\ \text{avec }t\geq r^{\pi_0},\ s\geq2
\tag*{(3.15.)}
\]
(\uline{\unsure{d'où|alors}} $s'\geq2^{(r^{\pi_0})}>r^{\pi_0}\geq r$), $\underline{\mathrm{Fl}}^{\,s'}M$ est stable sous le
foncteur $p_1^{C_0}:\underline{\mathrm{Fl}}\,M\unsure{\text{---}|\longrightarrow}\underline{\mathrm{Fl}}\,M$.

En effet, si on applique la première partie du
lemme \unsure{au} couple $(\unsure{r|s},s')$, on trouve que
$X_1(f)$ est dans $M^{s''}$, avec $s''={s'}^{(r^{\pi_0})}$. Mais
\[
s''={s'}^{(r^{\pi_0})}=(s^t)^{(r^{\pi_0})}
=s^{t.r^{\pi_0}}\underset{\substack{\text{car }t\geq r^{\pi_0}\\\text{donc }tr^{\pi_0}=t}}{=}
s^t=s'
\]
donc on voit que $f\in\mathrm{Fl}^{s'}$ implique $X_1(f)\in M^{s'}$, ou
ce qui qui revient au même, $p_1(f)\in\mathrm{Fl}^{s'}$.

\uline{Corollaire} Soient $C_0,r,s'=s^t$ ($t\geq r^{\pi_0}$, \inserted{$s\geq2$}) comme
dessus, et soit $\alpha$ un nombre ordinal \deleted{de} \inserted{tel}
\deleted{cardinal} que
\[
\operatorname{card}\alpha\leq s'=s^t.
\]
Alors $\underline{\mathrm{Fl}}^{\,s'}M$ est stable par le foncteur $p_\alpha=p_\alpha^{C_0}$.

On le voit aussitôt par récurrence transfinie,
\unsure{compte} \deleted{\unsure{donc}} \unsure{tenu} du lemme 3.3, et de l'hyp. de
stabilité 3.13.

Pour résumer

\uline{Théorème 2} \unsure{Soit|Soit,} $M$ une catégorie stable par \inserted{petites} $\underrightarrow{\lim}$
filtrantes, et munie d'une filtration cardinale
$(\mathrm{Filt}^\pi)_{\pi\geq\pi_0}$, telle que les $\mathrm{Filt}^\pi$ soient stables
par $\underrightarrow{\lim}$ filtrantes de cardinal $\leq\pi$ (pas forcément
grandes devant $\pi_0$), et par \inserted{les} sommes amalgamées simples
qui existent dans $M$. On suppose que les éléments
de $\mathrm{Filt}^{\pi_0}$ soient accessibles, et on désigne par
$\pi_1$ le plus petit cardinal $\geq\pi_0$, tel que les objets de
$\mathrm{Filt}^{\pi_0}$ soient $\pi_1$-accessibles. On pose \deleted{de plus}
